\chapter{系统关键模块实现}

\section{系统实现框架}

本系统主要使用 Python 开发，围绕设备接入、视频解码、目标检测、目标记忆、PTZ 控制和人机交互等模块展开实现。系统依托海康威视 HCNetSDK、PlayCtrl 播放库、ultralytics、OpenCV、Pillow 和 NumPy 等工具链，构建从实时码流输入到云台控制输出的闭环流程。实现过程中将设备接口、视频采集、检测分析和目标记忆分别封装，使各模块能够相对独立地调试和验证。

系统采用 Python 的主要原因在于其具有较好的算法生态和工程验证效率。ultralytics 能够较方便地加载 YOLOv8n 模型，OpenCV 和 Pillow 能够完成图像读取、格式转换和可视化绘制，ctypes 则可以调用海康 SDK 提供的动态库接口。虽然 Python 在运行效率和长期稳定性方面不如 C/C++，但对于本科毕业设计阶段的原型验证而言，Python 能够明显降低开发成本，使课题重点集中在检测、控制和系统闭环设计上。

系统运行流程由初始化、感知分析、控制执行和资源释放四个阶段组成。初始化阶段完成模型加载、设备接口初始化和实时预览启动；感知分析阶段持续读取最新视频帧并更新人体检测结果；控制执行阶段根据目标偏差和当前系统状态生成 PTZ 指令；资源释放阶段停止云台动作并释放预览、播放库和设备连接资源。该流程的完整性对真实设备尤为重要：若初始化顺序不合理，可能导致预览端口无法获取；若退出时未停止云台或释放句柄，则可能影响下次连接和设备端状态。

\section{视频采集与解码模块实现}

视频采集模块基于海康 SDK 的实时预览接口实现。系统首先调用 NET\_DVR\_Init 初始化 SDK，并通过 NET\_DVR\_Login\_V30 完成设备登录认证。登录成功后，系统构造 NET\_DVR\_PREVIEWINFO 结构体并调用 NET\_DVR\_RealPlay\_V40 开启实时预览。在该过程中，设备登录返回的用户 ID 和实时预览返回的句柄是后续控制和资源释放的重要依据，因此程序需要在每一步检查返回值，并在失败时记录对应错误码。

实时预览回调 RealDataCallBack\_V30 根据数据类型处理码流。当收到系统头数据时，系统调用 PlayM4\_OpenStream 打开播放流，并设置解码回调 DecCBFun；当收到码流数据时，系统调用 PlayM4\_InputData 将数据送入播放库。解码回调在获得图像帧后，将最新帧保存为 BMP 或 JPG 文件，并更新 latest\_frame\_data。检测线程和界面预览函数通过 capture\_frame\_for\_yolo 读取该最新帧，并统一缩放为 $1024\times768$ 尺寸，以便后续检测和控制使用。

该实现方式的优势是接口清晰、调试方便，保存的中间图像可直接用于检查解码结果；不足是文件读写会带来额外开销。为降低帧积压风险，系统并不对所有帧建立队列，而是始终读取最新可用帧。这样即使检测线程偶尔变慢，控制模块仍倾向于使用最近画面，避免因处理旧帧导致 PTZ 控制滞后。

为兼顾画质和速度，系统提供 PREVIEW\_USE\_BMP 参数。BMP 模式图像损失较小，有利于检测和显示，但磁盘 I/O 开销较大；JPG 模式速度更快，但存在压缩损失。当前系统默认使用 BMP 模式以提升画面质量。在后续嵌入式部署阶段，可将帧数据直接从解码回调转换为内存中的 OpenCV 图像，减少磁盘访问，并进一步降低端到端延迟。

\section{YOLO 人体检测模块实现}

系统使用 ultralytics 接口加载 \texttt{yolov8n.pt} 模型。YOLO 检测线程在 yolo\_enabled 状态为真时运行，每隔约 0.02 s 尝试读取一帧图像。为适配实时跟踪需求，系统设置 YOLO\_IMGSZ 为 384、YOLO\_CONF 为 0.4，并开启 YOLO\_HALF 以便在支持 GPU 的环境下使用半精度推理。上述参数体现了原型系统对实时性的优先考虑：较小输入尺寸能够降低推理耗时，置信度阈值能够过滤低质量检测框，半精度推理则为 GPU 环境下的加速提供可能。

YOLO 推理结果中包含多个检测框，系统仅保留类别编号为 0 的人体目标，并过滤低于置信度阈值的结果。对每个人体目标，系统提取边界框 $(x_1,y_1,x_2,y_2)$、中心点、置信度和面积，同时调用 extract\_body\_feature 生成低维人体特征。候选人体列表随后输入 ReID 记忆模块，用于跨帧身份关联。保留面积信息的原因在于，自动模式下系统默认选择画面中面积最大的可见人体，这通常对应距离摄像机更近、现场关注度更高的目标。

检测结果反馈模块用于呈现人体框、目标中心点、目标 ID、置信度、画面中心点、目标偏移量和当前跟踪模式。被选中的跟踪目标与其他可见人体采用不同标记方式，以便分析多目标场景中的目标选择逻辑。该反馈模块不仅用于展示检测结果，也用于验证 PTZ 控制输入是否合理，例如目标中心是否计算正确、倒挂模式下偏差方向是否反转、指定 ID 是否被正确选中。

\section{人体 ReID 记忆模块实现}

\texttt{reid.py} 中定义了 \texttt{BodyReIDMemory} 类。该模块维护一个 tracks 字典，每个 track 保存目标 ID、人体特征、首次出现帧、最近出现帧、边界框、中心点、面积和置信度等信息。新一帧检测结果到来时，模块先筛选短时间内仍有效的历史轨迹，然后对每个检测目标计算其特征与历史轨迹特征的余弦相似度。若最高相似度超过阈值，则复用历史 ID；否则创建新的 person\_id。为避免同一帧中两个检测目标复用同一个历史 ID，模块使用 assigned\_ids 记录本帧已经分配的编号，从而保证一帧内的目标编号互不冲突。

人体特征由 make\_body\_feature 构造，包括人体框区域的 BGR 平均颜色、宽高比、归一化宽度和归一化高度。该特征维度低、计算成本小，适合原型系统中的短时身份关联。颜色特征能够粗略描述人体外观，几何特征能够反映目标在画面中的尺度变化，两者结合可以在不额外运行深度 ReID 网络的情况下提供基本区分能力。为缓解单帧颜色波动，系统在匹配成功后使用特征融合更新历史特征：

\begin{equation}
    \symbf{f}_{new} = (1-\beta)\symbf{f}_{old} + \beta\symbf{f}_{det}
\end{equation}

其中 $\beta$ 为特征融合系数。当前系统中 ReID 匹配阈值设置为 0.86，最大缺失帧数设置为 120，表示目标短时消失后仍可保留一段记忆。目标选择函数在自动模式下选择面积最大的目标，在指定 ID 模式下仅返回对应 ID 的目标；若指定 ID 暂时不可见，则返回空并等待其重新入镜。该策略使指定目标跟踪具有“记忆等待”能力，避免目标短暂被遮挡后系统立即切换到其他人体。

\section{自适应 PTZ 控制模块实现}

PTZ 控制由底层指令封装、目标居中调整和控制循环协同完成。底层控制函数调用 NET\_DVR\_PTZControl\_Other 发送开始转动指令，等待一个短脉冲时间后发送停止指令。为避免指令过密，系统通过最近控制时间和控制冷却参数限制最小控制间隔。该函数只负责执行单次方向动作，不直接判断目标位置，从而使底层设备控制与上层视觉决策保持相对独立。

目标居中调整函数负责根据目标中心点计算偏差。函数首先利用平滑系数对目标坐标进行指数平滑，随后计算目标相对画面中心的水平偏差和垂直偏差。当偏差超过阈值时，函数选择对应的水平或垂直控制方向。偏差绝对值之和用于调整单次脉冲时长，使偏差较大时云台动作更明显。由于 SDK 控制接口表现为方向开始和方向停止，本文采用短脉冲方式近似连续控制，既能避免云台长时间转动，也便于根据偏差大小调节控制强度。

布控球在倒挂安装时，画面方向与云台控制方向会发生反转。系统引入倒挂状态变量，对预览画面、检测画面和 PTZ 控制偏差进行统一修正，从而保持观察方向与控制方向一致。该功能虽然不是核心算法，但对真实部署十分重要，因为现场安装条件经常受支架、供电和视角限制，倒挂适配可以降低重新安装设备的成本。

\section{人机交互模块实现}

系统构建人机交互模块，用于提供视频预览、检测状态反馈、跟踪模式切换、目标选择和人工接管等功能。交互模块与检测线程和控制线程共享必要状态，使测试人员能够观察目标识别结果、跟踪目标编号和云台控制效果。与复杂业务系统不同，本文交互模块的设计目标是服务于闭环验证和现场调试，使系统状态能够被及时观察和干预。

交互模块通过状态变量控制检测、自动跟踪、目标选择和倒挂适配等功能。自动模式下，系统默认选择画面中面积最大的可见人体；指定目标模式下，系统根据当前目标 ID 维持对特定人体的跟踪；当自动跟踪未开启、目标丢失或参数调试阶段出现异常时，人工接管能力可以作为安全补充。通过这种设计，系统既保留了自动闭环控制能力，也具备必要的人机协同调试空间。

\section{异常处理与资源释放}

系统在多个关键环节加入异常处理。YOLO 模型加载失败时会关闭检测功能并提示错误；检测线程运行时捕获推理异常并继续下一轮检测；视频帧读取失败时返回空结果，避免控制模块使用无效数据；PTZ 控制失败时输出 SDK 错误码。退出时，on\_close 函数会先调用 ptz\_stop 停止云台运动，再释放预览和 SDK 资源，降低设备端残留状态对下次运行的影响。

异常处理的设计目标不是掩盖错误，而是让系统在可恢复场景下尽量继续运行，并在不可恢复场景下有序退出。例如，单帧检测失败通常不影响下一帧检测，不应导致程序终止；但设备登录失败或预览启动失败属于关键链路异常，程序需要停止后续流程并清理已初始化资源。通过区分不同异常的严重程度，系统能够在原型阶段获得更好的调试体验。
